\exercisesection{}

\begin{enumerate}
\def\labelenumi{\arabic{enumi})}
\tightlist
\item
  在 \(\mathbf{R}^2\) 中，设 \(\lambda\) 与 \(\mu\) 为正测度
\end{enumerate}

\[
f \mapsto \int_ {0} ^ {+ \infty} f (x, 0) d x, \quad f \mapsto \int_ {0} ^ {+ \infty} f (- x, x) d x \quad \left(f \in \mathcal {K} (\mathbf {R} ^ {2})\right).
\]

证明 \(\lambda\) 与 \(\mu\) 是可卷积的。设 \(u\) 是同态 \((x,y) \mapsto x\) （\(\mathbf{R}^2\) 到 \(\mathbf{R}\) 上的）。证明 \(u\) 是 \(\lambda\) -逆紧的且是 \(\mu\) -逆紧的，但 \(u(\lambda)\) 与 \(u(\mu)\) 不是可卷积的。

\begin{enumerate}
\def\labelenumi{\arabic{enumi})}
\setcounter{enumi}{1}
\item
  设 X 是一个局部紧空间，一个局部紧群 G 在其中连续地从左边作用。设 E 是 \(\mathcal{M}(X)\) 的一个线性子空间，在诸 \(\boldsymbol{\gamma}(s)\) （ \(s \in G\) ）下稳定，赋有一个比 \(\mathcal{K}(X)\) 中紧收敛拓扑更细的拟完备局部凸拓扑。对每个 \(s \in G\) ，设 \(\boldsymbol{\gamma}_{\mathrm{E}}(s)\) 为 \(\boldsymbol{\gamma}(s)\) 在 E 上的限制。假设 \(\mu \in E\) 蕴含 \(|\mu| \in E\) ，且 G 在 E 中的表示 \(\gamma_{E}\) 是等度连续的。那么，若 \(\mu \in \mathbf{E}\) 且 \(\nu \in \mathcal{M}^1 (\mathbf{G})\) ，则 \(\nu\) 与 \(\mu\) 可卷积，并且 \(\nu * \mu = \gamma_{\mathbf{E}}(\nu)\mu \in \mathbf{E}\) 。（特别利用 Ch. VI, §1, No.~7 的 Prop. 17。）
\item
  对 \(x = (x_{1}, \ldots, x_{n}) \in \mathbf{R}^{n}\) ，置 \(|x|^{2} = x_{1}^{2} + \cdots + x_{n}^{2}\) 。设 \(M_{1}\) 为这样的测度 \(\mu\) （在 \(R^{n}\) 上）所成的集合：存在一个实数 k 使得函数 \((1 + |x|^{2})^{k}\) 是 \(\mu\) -可积的。设 \(M_{2}\) 为这样的测度 \(\nu\) （在 \(R^{n}\) 上）所成的集合：对每个 k，函数 \((1 + |x|^{2})^{k}\) 都是 \(\nu\) -可积的。证明：若 \(\mu \in M_{1}\) 且 \(\nu \in M_{2}\) ，则 \(\mu\) 与 \(\nu\) 可卷积，且 \(\mu * \nu \in M_{1}\) ；若 \(\mu \in M_{2}\) 且 \(\nu \in M_{2}\) ，则 \(\mu * \nu \in M_{2}\) 。（先证明：若 \(u \geqslant 0\) ，\(v \geqslant 0\) ，则
\end{enumerate}

\[
(1 + u ^ {2}) (1 + v ^ {2}) \geqslant \frac {1}{3} (1 + (u + v) ^ {2});
\]

由此推出，对任意的 \(x, y\) （在 \(\mathbf{R}^n\) 中），

\[
1 + | x | ^ {2} \leqslant 3 (1 + | y | ^ {2}) (1 + | x + y | ^ {2}).
\]

然后设 \(\mu \in \mathcal{M}_1\) ，\(\nu \in \mathcal{M}_2\) ，其中 \(\mu \geqslant 0\) ，\(\nu \geqslant 0\) 。设 \(f\) 是一个连续函数 \(\geqslant 0\) （在 \(\mathbf{R}^n\) 上）。存在一个 \(k\) 使得 \(\mu = (1 + |x|^2)^k \cdot \mu_1\) ，其中 \(\mu_1\) 有界；令 \(\nu_1 = (1 + |x|^2)^k \cdot \nu\) ，它是有界的。于是

\[
\int^ {*} f (x + y) d \mu (x) d \nu (y) \leqslant 3 ^ {k} \int^ {*} (1 + | x | ^ {2}) ^ {k} f (x) d (\mu_ {1} * \nu_ {1}) (x).
\]

由此得到 \(\mu\) 与 \(\nu\) 的可卷积性以及 \(\mu * \nu \in \mathcal{M}_1\) 这一事实。对 \(\mu, \nu\) （\(\mathcal{M}_2\) 中的）类似地论证。）

\begin{enumerate}
\def\labelenumi{\arabic{enumi})}
\setcounter{enumi}{3}
\item
  设 \(\mu\) 是 \(\mathbf{R}\) 上的 Lebesgue 测度，\(\nu\) 是 \([0, +\infty[\) 上的 Lebesgue 测度。设 \(x_1, x_2\) 属于 \(\mathbf{R}\) 。证明卷积 \(((\varepsilon_{x_1} - \varepsilon_{x_2}) * \nu) * \mu\) 有定义，但 \(\mu\) 与 \(\nu\) 不可卷积。证明卷积 \(\nu * ((\varepsilon_{x_1} - \varepsilon_{x_2}) * \mu)\) 与 \((\nu * (\varepsilon_{x_1} - \varepsilon_{x_2})) * \mu\) 都有定义，但当 \(x_1 \neq x_2\) 时它们不相等。
\item
  设 G 是紧群，\(\mu\) 是 G 上一个以 G 为支集的正测度，使得 \(\mu*\mu=\mu\) 。证明 \(\mu\) 是 G 的规范化 Haar 测度。（先证明 \(\|\mu\|=1\) 。然后，若 \(\mu\) 不是 G 的 Haar 测度，则存在一个函数 \(f\in\mathcal{K}_{+}(G)\) 使得
\end{enumerate}

\[
\int f (s) d \mu (s) \geqslant \int f (s t) d \mu (s)
\]

对所有 \(t \in \mathbf{G}\) 成立，且 \(\int f(s) d\mu(s) > \int f(st) d\mu(s)\) 对某个 \(t\) 成立。证明这时有 \((\mu * \mu)(f) > \mu(f)\) 。

\begin{enumerate}
\def\labelenumi{\arabic{enumi})}
\setcounter{enumi}{5}
\tightlist
\item
  \begin{enumerate}
  \def\labelenumii{\alph{enumii})}
  \tightlist
  \item
    设 \(I = [0,1]\) 。设 \(f\) 是一个连续函数 \(\geqslant 0\) （在 \(\mathbf{R}\) 上），其支集包含在 \([-1,0]\) 中。证明函数 \(\pmb{\gamma}(s)f|I\) （ \(s \in I\) ）所成的集合在 \(\mathcal{K}(I)\) 中有无穷秩。
  \end{enumerate}
\end{enumerate}

\begin{enumerate}
\def\labelenumi{\alph{enumi})}
\setcounter{enumi}{1}
\item
  设 \(f_1, \ldots, f_n\) 属于 \(\mathcal{K}(\mathbf{R})\) 。设 \(M\) 为这样的测度 \(\mu \in \mathcal{M}(I)\) 所成的集合：\(\mu(f_1) = \ldots = \mu(f_n) = 0\) 。证明函数 \(y \mapsto \int f(x + y) d\mu(x)\) （ \(y \in I\) ）（其中 \(\mu\) 取遍 \(M\) ）所成的集合在 \(\mathcal{K}(I)\) 中有无穷秩。（利用 \(a\) ）。
\item
  设 \(g_1, \ldots, g_p\) 属于 \(\mathcal{K}(\mathbf{R})\) 。由 \(b\) 推出，存在 \(\mu \in \mathbb{M}\) 与 \(\nu \in \mathcal{M}(\mathrm{I})\) 使得 \(\nu(g_1) = \ldots = \nu(g_p) = 0\) ，\((\nu * \mu)(f) \neq 0\) 。
\item
  由 c) 推出，映射 \((\mu,\nu)\mapsto\nu*\mu\) （\(\mathcal{M}(\mathbf{T})\times\mathcal{M}(\mathbf{T})\) 到 \(\mathcal{M}(\mathbf{T})\) 中的）不是淡连续的。
\end{enumerate}

\begin{enumerate}
\def\labelenumi{\arabic{enumi})}
\setcounter{enumi}{6}
\tightlist
\item
  设 \(G\) 为局部紧群。
\end{enumerate}

\begin{enumerate}
\def\labelenumi{\alph{enumi})}
\item
  证明：若正测度 \(\mu \in \mathcal{C}'(\mathrm{G})\) 在 \(\mathcal{C}'(\mathrm{G})\) 中有一个作为正测度的逆，则 \(\mu\) 必是点测度。
\item
  取 G 为有限群 \(\mathbf{Z}/2\mathbf{Z}\) ，给出一个非点测度 \(\mu \in \mathcal{M}(\mathbf{G})\) 的例子，它在 \(\mathcal{M}(\mathbf{G})\) 中可逆。
\end{enumerate}

\begin{enumerate}
\def\labelenumi{\arabic{enumi})}
\setcounter{enumi}{7}
\item
  设 G 为局部紧群；用 \(\mathcal{C}_{+}^{\prime}(G)\) 表示 G 上具有紧支集的正测度所成的集合。证明映射 \((\mu, \nu) \mapsto \mu * \nu\) （\(\mathcal{M}_{+}(G) \times \mathcal{C}_{+}^{\prime}(G)\) 到 \(\mathcal{M}_{+}(G)\) 中的）是连续的，其中 \(\mathcal{C}'(G)\) 赋有弱拓扑 \(\sigma(\mathcal{C}'(G), \mathcal{C}(G))\) 而 \(\mathcal{M}(G)\) 赋有淡拓扑 \(\sigma(\mathcal{M}(G), \mathcal{K}(G))\) 。（利用 Ch. III, §1 的 Exer. 10 以及 G 是仿紧的这一事实。）
\item
  \begin{enumerate}
  \def\labelenumii{\alph{enumii})}
  \tightlist
  \item
    设 G 为局部紧群，B 为 \(\mathcal{M}(\mathbf{G})\) 的一个有界子集，C 为 \(\mathcal{C}'(\mathbf{G})\) 的一个等度连续子集；证明：若 \(\mathcal{C}'(\mathbf{G})\) 赋有弱拓扑 \(\sigma(\mathcal{C}'(\mathbf{G}), \mathcal{C}(\mathbf{G}))\) ，且 \(\mathcal{M}(\mathbf{G})\) 赋有淡拓扑 \(\sigma(\mathcal{M}(\mathbf{G}), \mathcal{K}(\mathbf{G}))\) ，则映射 \((\mu, \nu) \mapsto \mu * \nu\) （B × C 到 \(\mathcal{M}(\mathbf{G})\) 中的）是连续的。（注意所有测度 \(\nu \in \mathbf{C}\) 的支集都在一个公共紧集内，并利用 Ch. III, §4, No.~3 中 Prop. 6 后的注。）
  \end{enumerate}
\end{enumerate}

\begin{enumerate}
\def\labelenumi{\alph{enumi})}
\setcounter{enumi}{1}
\tightlist
\item
  若在 \(\mathcal{M}^1 (\mathbf{R})\) 中令 \(\mu_n = \varepsilon_n\) ，\(\nu_{n} = \varepsilon_{-n}\) ，则序列 \((\mu_n)\) 与 \((\nu_{n})\) 对于弱拓扑 \(\sigma (\mathcal{M}^1 (\mathbf{R}),\overline{\mathcal{K}(\mathbf{R})})\) 淡地趋于 0，但序列 \((\mu_n*\nu_n)\) 对于淡拓扑不趋于 0。
\end{enumerate}

¶ 10) 对局部紧空间 T，用 \(\mathcal{C}^{\infty}(T)\) 表示 T 上有界连续数值函数所成的 Banach 空间。称 \(\mathcal{M}^{1}(T)\) 的子集 H 是挤压的，如果对每个 \(\varepsilon > 0\) ，存在 T 的一个紧子集 K，使得 \(|\mu|(T - K) \leqslant \varepsilon\) 对所有 \(\mu \in H\) 都成立。

\begin{enumerate}
\def\labelenumi{\alph{enumi})}
\item
  证明：若 \(\mathcal{M}^{1}(G)\) 的子集 H 是挤压的，且对 \(\mathcal{M}^{1}(T)\) 的范数所定义的拓扑有界，则它对拓扑 \(\sigma(\mathcal{M}^{1}(T),\mathcal{C}^{\infty}(T))\) 是相对紧的（注意 H 对 \(\sigma(\mathcal{M}^{1}(T),\mathcal{K}(T))\) 是相对紧的）。
\item
  再假设 T 是仿紧的，反之证明：若 H 是 \(\mathcal{M}^1 (\mathrm{T})\) 中对 \(\sigma (\mathcal{M}^1 (\mathrm{T}),\mathcal{C}^\infty (\mathrm{T}))\) 相对紧的子集，则 H 对 \(\mathcal{M}^1 (\mathrm{T})\) 的范数有界且是挤压的。（先考虑 \(\mathbf{T} = \mathbf{N}\) 的情形，应用 Ch. V, §5 的 Exer. 15。然后考察 T 在无穷处可数、是一列相对紧开集 \((\mathrm{U}_n)\) （满足 \(\overline{\mathrm{U}}_n\subset \mathrm{U}_{n + 1}\) ）之并的情形。用反证法证明，可以化归为这样的情形：对每个 \(n\) 存在一个定义在 T 上的连续数值函数 \(f_{n}\) ，其支集包含在 \(\mathrm{U}_{n + 1} - \overline{\mathrm{U}}_n\) 中，满足 \(\| f_n\| \leqslant 1\) ，并且存在一列属于 H 的测度 \((\mu_n)\) ，使得 \(\mu_{n}(f_{n})\geqslant a > 0\) 对所有 \(n\) 成立。于是考虑连续映射 \(u:\mathrm{L}^{1}(\mathbf{N})\to \mathcal{C}^{\infty}(\mathrm{T})\)
\end{enumerate}

使得 \(u((\xi_n)) = \sum_{n=0}^{\infty} \xi_n f_n\) ，并与对 \(T = N\) 已证明的结果相矛盾。

通过考虑 \(u\) 的转置。最后，当 \(T\) 是任意的仿紧局部紧空间时，\(T\) 是一族在无穷处可数的局部紧空间 \((T_{\alpha})\) 的拓扑和；对每个 \(\alpha\) ，令 \(m_{\alpha} = \sup_{\mu \in H} |\mu|(T_{\alpha})\) ；用反证法并利用前一种情形证明，必然有 \(m_{\alpha} = 0\) ，除对可数无穷多个指标 \(\alpha\) 外。）

\begin{enumerate}
\def\labelenumi{\alph{enumi})}
\setcounter{enumi}{2}
\tightlist
\item
  证明 \(b\) ) 的结论对 Ch. IV, §4 的 Exer. 16 h) 中定义的非仿紧局部紧空间不成立。
\end{enumerate}

¶ 11) 设 G 为局部紧群；在 \(\mathcal{M}^{1}(G)\) 上，用 \(\mathcal{T}_{\mathrm{I}}\) 表示拓扑 \(\sigma(\mathcal{M}^{1}(G), \overline{\mathcal{K}(G)})\) ，用 \(\mathcal{T}_{\mathrm{III}}\) 表示拓扑 \(\sigma(\mathcal{M}^{1}(G), \mathcal{C}^{\infty}(G))\) ，并用 \(\mathcal{M}_{\mathrm{I}}\) 、\(\mathcal{M}_{\mathrm{III}}\) 分别表示空间 \(\mathcal{M}^{1}(G)\) 赋有 \(\mathcal{T}_{\mathrm{I}}\) 与 \(\mathcal{T}_{\mathrm{III}}\) 后所得的空间。

\begin{enumerate}
\def\labelenumi{\alph{enumi})}
\item
  设 A 是 \(\mathcal{M}_{\mathrm{I}}\) 的一个有界子集，B 是 \(\mathcal{M}_{\mathrm{III}}\) 的一个相对紧子集。证明 \(\mathrm{A} \times \mathrm{B}\) 上映射 \((\mu, \nu) \mapsto \mu * \nu\) （\(\mathcal{M}_{\mathrm{I}} \times \mathcal{M}_{\mathrm{III}}\) 到 \(\mathcal{M}_{\mathrm{I}}\) 中的）的限制是连续的（利用 Exer. 10，以化归为估计积分 \(\iint f(st) d\mu(s) d\nu(t)\) ，当 \(f(st) = \sum_{i} u_{i}(s)v_{i}(t)\) 、其中 \(u_{i}\) 与 \(v_{i}\) 属于 \(\mathcal{K}(\mathrm{G})\) 时）。
\item
  给出一个 G 为紧（从而 \(T_{I} = T_{III}\) ）的例子，表明 \((\mu, \nu) \mapsto \mu * \nu\) 作为 \(M_{I} \times M_{III}\) 到 \(M_{I}\) 中的映射，既不关于 \(M_{I}\) 的相对紧子集亚连续，也不关于 \(M_{III}\) 的相对紧子集亚连续 (cf.~Exer. 6)。
\item
  设 \(A, B\) 为 \(\mathcal{M}_{\mathrm{III}}\) 的两个相对紧子集。证明 \(A \times B\) 上映射 \((\mu, \nu) \mapsto \mu * \nu\) （\(\mathcal{M}_{\mathrm{III}} \times \mathcal{M}_{\mathrm{III}}\) 到 \(\mathcal{M}_{\mathrm{III}}\) 中的）的限制是连续的（方法同 a)）。
\item
  用 E 表示由 G 的开子集的特征函数的线性组合构成的 \(\mathbf{R}^{\mathrm{G}}\) 的子空间，用 \(\mathcal{T}_{\mathrm{IV}}\) 表示拓扑 \(\sigma(\mathcal{M}^{1}(\mathrm{G}),\mathrm{E})\) ，并用 \(\mathcal{M}_{\mathrm{IV}}\) 表示空间 \(\mathcal{M}^{1}(\mathrm{G})\) 赋有 \(\mathcal{T}_{\mathrm{IV}}\) 后所得的空间。回忆一下，\(\mathcal{T}_{\mathrm{III}}\) 与 \(\mathcal{T}_{\mathrm{IV}}\) 在由正测度组成的 \(\mathcal{M}^{1}(\mathrm{G})\) 的有界子集上诱导的拓扑一般是不相同的 (Ch. V, §5, Exer. 16 c))。再回忆一下，\(\mathcal{M}_{\mathrm{IV}}\) 的紧子集与 Banach 空间 \(\mathcal{M}^{1}(\mathrm{G})\) 赋弱化拓扑 \(\sigma(\mathcal{M}^{1}(\mathrm{G}),(\mathcal{M}^{1}(\mathrm{G}))'\) 时 \(\mathcal{M}^{1}(\mathrm{G})\) 的紧子集相同 (Ch. VI, §2, Exer. 12)。证明：若 A、B 是 \(\mathcal{M}_{\mathrm{IV}}\) 的两个相对紧子集，则映射 \((\mu,\nu)\mapsto\mu*\nu\) （\(\mathcal{M}_{\mathrm{IV}}\times\mathcal{M}_{\mathrm{IV}}\) 到 \(\mathcal{M}_{\mathrm{IV}}\) 中的）在 A × B 上的限制是连续的。（只限于 A 与 B 是紧的情形，并且先证明这时 A × B 在上述映射下的像对 \(\sigma(\mathcal{M}^{1}(\mathrm{G}),(\mathcal{M}^{1}(\mathrm{G}))'\) 是紧的；为此，应用 Eberlein 定理与 Šmulian 定理 (TVS, IV, §5, No.~3)，从而化归为证明：若 \((\mu_{n})\) 、\((\nu_{n})\) 是在 \(\mathcal{M}_{\mathrm{IV}}\) 中收敛于 0 的两个序列，则序列 \((\mu_{n}*\nu_{n})\) 亦然；利用 Ch. V, §5 的 Prop. 12 与 Exer. 15。最后，为证 \((\mu,\nu)\mapsto\mu*\nu\) 对 \(\mathcal{T}_{\mathrm{IV}}\) 连续，利用 c) 以及 \(\mathcal{T}_{\mathrm{III}}\) 比 \(\mathcal{T}_{\mathrm{IV}}\) 粗这一事实。）
\item
  取 \(\mathbf{G} = \mathbf{R}^2\) ；设 \(a, b\) 是 \(\mathbf{G}\) 在 \(\mathbf{R}\) 上的标准基的向量，\(\rho_n\) 是区间 \(\mathbf{I} = [0, \pi]\) （\(\mathbf{R}\) 中的）上以函数 \(\sin(2^n x)\) 为密度（关于 Lebesgue 测度）的测度，\(\mu_n\) 是测度 \(\rho_n \otimes \varepsilon_0\) （\(\mathbf{R} \times \mathbf{R}\) 上的）；另一方面，设 \(\nu_n = \varepsilon_{b/2^n} - \varepsilon_0\) （\(\mathbf{G}\) 上的）。证明序列 \((\mu_n)\) 在 \(\mathcal{M}_{\mathrm{IV}}\) 中趋于 0，序列 \((\nu_n)\) 在 \(\mathcal{M}_{\mathrm{III}}\) 中趋于 0，但序列 \((\mu_n * \nu_n)\) 在 \(\mathcal{M}_{\mathrm{IV}}\) 中不趋于 0 (cf.~Ch. V, §5, Exer. 16)。
\item
  给出一个 G 为紧且 \((\mu,\nu)\mapsto\mu*\nu\) 作为 \(M_{IV}\times M_{IV}\) 到 \(M_{IV}\) 中的映射不关于 \(M_{IV}\) 的紧子集亚连续的例子（方法同 Exer. 6，注意对给定的 \(f\in E\) ，存在紧子集 \(H\subset M_{IV}\) 使得函数 \(t\mapsto\int f(st)d\mu(s)\) （其中 \(\mu\) 取遍 H）所成的集合在 R 上有有限秩）。
\end{enumerate}

\begin{enumerate}
\def\labelenumi{\arabic{enumi})}
\setcounter{enumi}{11}
\tightlist
\item
  设 \(G\) 为一个非幺模的局部紧群。
\end{enumerate}

\begin{enumerate}
\def\labelenumi{\alph{enumi})}
\item
  证明存在一个有界正测度 \(\mu\) （\(\mathbf{G}\) 上的），使得 \(\Delta_{\mathbf{G}} \cdot \mu\) 不是有界的（取 \(\mu\) 为离散的）。
\item
  设 \(\mu'\) 是 G 上的一个左 Haar 测度。那么 \(\mu\) 与 \(\mu'\) 可卷积 (Prop. 5)。证明 \(\mu'\) 与 \(\mu\) 不可卷积。
\end{enumerate}

\begin{enumerate}
\def\labelenumi{\arabic{enumi})}
\setcounter{enumi}{12}
\tightlist
\item
  设 \(r\) 是一个满足 \(0 < r < 1\) 的数；对每个整数 \(n \geqslant 1\) ，用 \(\lambda_{n,r}\) 表示测度 \((\varepsilon_{r^n} + \varepsilon_{-r^n}) / 2\) （\(\mathbf{R}\) 上的），并令 \(\mu_{n,r} = \lambda_{1,r} * \lambda_{2,r} * \dots * \lambda_{n,r}\) 。
\end{enumerate}

\begin{enumerate}
\def\labelenumi{\alph{enumi})}
\item
  证明序列 \((\mu_{n,r})\) 淡收敛到 R 上的一个测度 \(\mu_{r}\) ，其支集包含在 I = {[}-1,+1{]} 中（证明对 R 的每个区间 U，序列 \((\mu_{n,r}(U))\) 收敛）。
\item
  证明当 \(r < 1/2\) 时，测度 \(\mu_r\) 与 \(\mathbf{R}\) 上的 Lebesgue 测度互奇异，但 \(\mu_{1/2}\) 是 Lebesgue 测度在 I 上诱导的测度。
\item
  设 \(\nu_{1/4}\) 是 \(\mu_{1/4}\) 在位似 \(t \mapsto 2t\) （\(\mathbf{R}\) 中的）下的像。证明 \(\mu_{1/4} * \nu_{1/4} = \mu_{1/2}\) ，尽管 \(\mu_{1/4}\) 与 \(\nu_{1/4}\) 都与 Lebesgue 测度互奇异（利用 Exer. 11 a)）。
\end{enumerate}

